\documentclass[review,12pt,authoryear]{elsarticle}

\usepackage{amsmath,amssymb,mathtools}
\usepackage{booktabs}
\usepackage{enumitem}
\usepackage{microtype}
\journal{Advances in Space Research}

\newcommand{\E}{\mathbb{E}}
\newcommand{\dv}{\Delta v}
\newcommand{\state}{\mathcal{S}}
\newcommand{\scenset}{\Omega}

\begin{document}
\begin{frontmatter}

\title{A Rolling-Horizon Framework for Service-Level Agreement Acceptance and
Infrastructure Planning in On-Orbit Servicing}

\author[inst1]{Author Name\corref{cor1}}
\ead{author@example.com}
\cortext[cor1]{Corresponding author.}
\address[inst1]{Institution, Department, City, Country}

\begin{abstract}
On-orbit servicing (OOS) operators interact with potential customers through
service-level agreements (SLAs). Deciding whether to accept an SLA requires
more than testing the proposed contract in isolation: the operator must account
for its current fleet, depot inventory, previously accepted commitments,
uncertain launch opportunities, and the operational consequences of serving the
new customer. This paper presents a rolling-horizon decision framework that
links SLA acceptance to launch-capacity bookings, infrastructure deployment,
inventory replenishment, and service-vehicle routing. Candidate architectures
are generated through multi-objective search and evaluated using a
mixed-integer launch-manifest model under sampled launch and demand outcomes.
The resulting reliability measure is the probability that the complete pool of
accepted SLAs can be fulfilled, rather than a user-specified input. The
framework returns an accept, defer, reject, or accept-with-expansion
recommendation together with expected commercial value and operationally
feasible counterfactuals. A research demonstrator illustrates how a contract
can be feasible with existing assets, feasible only after resupply or
deployment, or commercially unattractive once uncertainty and prior
commitments are included. The formulation provides an explainable bridge
between customer-facing contract decisions and the physical planning of an OOS
network.
\end{abstract}

\begin{keyword}
on-orbit servicing \sep service-level agreement \sep rolling-horizon planning
\sep launch uncertainty \sep space logistics \sep multi-objective optimisation
\end{keyword}
\end{frontmatter}

\section{Introduction}
\label{sec:introduction}

On-orbit servicing is shifting from isolated demonstrations toward repeated
commercial operations involving inspection, refuelling, repair, relocation,
and life extension. An operator offering these services must decide both which
customer contracts to accept and what infrastructure to deploy and operate.
Treating these decisions separately can produce contracts that appear
profitable but cannot be delivered with the available vehicles, inventory, or
launch opportunities.

We represent customer commitments as SLAs. An SLA specifies a set of client
spacecraft, required services and quantities, service windows, contract value,
and a maximum acceptable loss boundary. The loss boundary limits the fraction
of client-asset value that may remain unserved over the contract horizon. The
operator receives candidate SLAs sequentially while the operational state and
information about future launches evolve. SLA acceptance is therefore a
rolling-horizon decision problem.

This paper makes three contributions. First, it formulates SLA acceptance as an
integrated commercial and operational decision rather than a static contract
screen. Second, it couples multi-objective service-architecture search to an
explicit launch-manifest and infrastructure-selection problem under uncertain
outcomes. Third, it retains both Pareto-optimal and dominated feasible
solutions, allowing rejected proposals to be accompanied by operationally
supported counterfactuals such as a later deadline, smaller scope, or additional
deployment.

\section{Rolling-horizon SLA decision problem}
\label{sec:problem}

At decision epoch $\tau$, the operator state is
\begin{equation}
 \state_\tau=(\mathcal{A}_\tau,\mathcal{V}_\tau,\mathcal{D}_\tau,
 \mathbf{q}_\tau,\mathcal{L}_\tau,\mathcal{M}_\tau,\mathcal{I}_\tau),
\end{equation}
where $\mathcal{A}_\tau$ is the pool of accepted SLAs,
$\mathcal{V}_\tau$ the available service vehicles,
$\mathcal{D}_\tau$ the deployed depots, $\mathbf{q}_\tau$ their inventories,
$\mathcal{L}_\tau$ the uncertain launch calendar, $\mathcal{M}_\tau$ the booked
manifests, and $\mathcal{I}_\tau$ the information available at epoch $\tau$.
All durations are stored in seconds and each state has an explicit timestamp.

A proposed SLA $k$ is
\begin{equation}
 c_k=(\mathcal{T}_k,\mathbf{d}_k,[t_k^-,t_k^+],R_k,\ell_k,P_k),
\end{equation}
where $\mathcal{T}_k$ is the target set, $\mathbf{d}_k$ contains service
quantities, $[t_k^-,t_k^+]$ is the service window, $R_k$ is contract revenue,
$\ell_k$ is the acceptable loss fraction, and $P_k$ specifies penalties. The
candidate is evaluated together with all earlier commitments:
\begin{equation}
 \mathcal{A}_\tau^+=\mathcal{A}_\tau\cup\{c_k\}.
\end{equation}

Let $\omega\in\scenset$ denote a joint scenario with probability $p_\omega$.
Each scenario contains launch time, destination, purchasable payload capacity,
delay or cancellation, service demand, and operational availability. These
variables are sampled jointly because launch date, destination, provider,
vehicle, and capacity are coupled. The optimisation accepts any weighted
scenario set, including samples from an empirical Bayesian launch model.

The framework returns: accept using existing assets; accept with a specified
resupply or deployment; defer to a later window; or reject.

\section{Integrated computational framework}
\label{sec:framework}

At every decision epoch, the framework:
\begin{enumerate}
 \item parses and validates the proposed SLA;
 \item combines it with previously accepted commitments;
 \item calculates time-dependent transfer costs;
 \item generates service routes and single-depot architectures;
 \item constructs multi-depot launch and deployment manifests;
 \item evaluates every plan across the joint scenario set;
 \item retains a feasible archive and extracts its Pareto front;
 \item recommends an action and reports counterfactual alternatives; and
 \item updates the rolling state after explicit acceptance.
\end{enumerate}

Natural-language processing only translates customer input into the contract
schema. Deterministic validation remains authoritative: missing targets,
ambiguous per-spacecraft quantities, invalid orbital elements, and absent
commercial terms trigger clarification before optimisation.

\section{Time-dependent transfer-cost model}
\label{sec:transfer}

The demonstrator uses a thrust--coast--thrust approximation based on Edelbaum's
analytical low-thrust transfer model \citep{edelbaum1961}. For circular-orbit
velocities $v_0$ and $v_f$, with plane change $\Delta i$,
\begin{equation}
 \dv_{\mathrm{E}}=
 \sqrt{v_0^2+v_f^2-2v_0v_f\cos\left(\frac{\pi}{2}\Delta i\right)}.
\end{equation}
The propellant requirement and powered duration are
\begin{align}
 m_p&=m_0\left[1-\exp\left(-\frac{\dv_{\mathrm{E}}}
 {I_{\mathrm{sp}}g_0}\right)\right],\\
 t_b&=\frac{m_p I_{\mathrm{sp}}g_0}{T}.
\end{align}

RAAN is closed using natural $J_2$ drift during coast arcs:
\begin{equation}
 \dot{\Omega}_{J_2}=-\frac{3}{2}J_2n
 \left(\frac{R_E}{a(1-e^2)}\right)^2\cos i.
\end{equation}
Departure epoch and time of flight therefore affect feasibility and cost.
Values are precomputed in a lookup table indexed by origin, destination,
departure epoch, and time of flight. The demonstrator uses 30-day departure
increments and flight times of 30, 60, 90, 120, 180, and 240 days. This
conservative model supports rapid architecture screening; high-fidelity
trajectory optimisation is required before execution.

\section{Multi-shuttle service routing}
\label{sec:routing}

For a depot and its assigned requests, multi-directional local search (MDLS)
explores service sequences and vehicle allocations \citep{tricoire2012}. A
routing solution $r$ is evaluated through
\begin{align}
 f_1(r)&=\sum_{v\in\mathcal{V}}\sum_{(i,j)\in r_v}\dv_{ij},\\
 f_2(r)&=\sum_{v\in\mathcal{V}}u_v,\\
 f_3(r)&=\sum_{k\in\mathcal{A}_\tau^+}W_k[1-y_k(r)],
\end{align}
where $u_v$ indicates vehicle use, $y_k(r)$ indicates fulfilment of SLA $k$,
and $W_k$ is protected asset value. These objectives represent operational
effort, fleet requirement, and unrecovered value.

The search produces non-dominated single-depot, multi-shuttle concepts for each
candidate depot. The full feasible archive is retained. Dominated plans are
excluded from the primary decision but remain useful for counterfactuals.

\section{Launch manifest and infrastructure selection}
\label{sec:manifest}

The architecture layer chooses candidate servicers, depots, and
commodity-specific resupply lots. Let $z_j$ indicate selection of item $j$ and
$x_{jl}^{\omega}$ its assignment to launch $l$ under scenario $\omega$. With
scenario capacity $C_l^\omega$ and compatibility $a_{jl}^\omega$,
\begin{align}
 \sum_lx_{jl}^\omega&=z_j &&\forall j,\omega,\\
 \sum_jm_jx_{jl}^\omega&\le C_l^\omega &&\forall l,\omega,\\
 x_{jl}^\omega&\le a_{jl}^\omega &&\forall j,l,\omega.
\end{align}
Prerequisite constraints ensure, for example, that resupply is selected only
for an existing or simultaneously deployed depot. For at most 16 optional
items, the demonstrator enumerates the binary manifest space exactly. A MILP
solver can replace enumeration at larger scale without changing the interface.

\section{Reliability and commercial decision model}
\label{sec:decision}

Let $y_{k\omega}(x)=1$ when architecture $x$ fulfils SLA $k$ in scenario
$\omega$. Its reliability is
\begin{equation}
 \rho_k(x)=\sum_{\omega\in\scenset}p_\omega y_{k\omega}(x).
\end{equation}
Launch disruption is embedded in $\rho_k$ because delayed, cancelled,
misdirected, or capacity-constrained launches alter fulfilment. Reliability is
therefore an output of the coupled model. The SLA loss boundary imposes
\begin{equation}
 \rho_k(x)\ge 1-\ell_k.
\end{equation}

Expected net value is
\begin{equation}
 \mathrm{ENV}(x)=
 \sum_{k\in\mathcal{A}_\tau^+}R_k\rho_k(x)
 -C_{\mathrm{launch}}(x)-C_{\mathrm{manufacture}}(x)
 -\E_\omega[P(x,\omega)]-c_{\dv}f_1(x).
\end{equation}
Among reliable and physically feasible plans, the recommended plan maximises
expected net value. An expansion decision is returned when new assets or
resupply are required. Deferral is returned when a later window yields a
feasible positive-value plan. Otherwise the proposal is rejected.

\section{State transition and counterfactuals}
\label{sec:rolling}

After acceptance, the selected manifest and reservations become commitments:
\begin{equation}
 \state_{\tau+1}=
 \Phi(\state_\tau,c_k,x_\tau,\xi_{\tau:\tau+1}),
\end{equation}
where $\xi_{\tau:\tau+1}$ contains observed launches, service events, demand,
and inventory changes. The next SLA is evaluated against the updated state.

For a rejected or marginal proposal, the framework searches the archive for
nearby feasible contracts. A distance measure may combine relative changes in
quantity, deadline, target scope, and revenue. The nearest alternatives are
Pareto filtered before presentation, so every counterfactual has an evaluated
operational plan behind it.

\section{Research demonstrator}
\label{sec:demonstrator}

The demonstrator exposes the framework through a local web interface and a
notebook. A user describes a contract in natural language; the system requests
missing or ambiguous information and displays the structured interpretation.
Results include the recommendation, reliability, expected value, launch
manifest, depot and vehicle use, service sequence, and feasible alternatives.

\begin{table}[ht]
\centering
\caption{Illustrative behaviour of the integrated decision framework.}
\label{tab:cases}
\begin{tabular}{p{0.25\linewidth}p{0.25\linewidth}p{0.38\linewidth}}
\toprule
Request & Recommendation & Operational interpretation\\
\midrule
20 kg xenon & Accept & Existing inventory and vehicle capacity cover the SLA.\\
100 kg xenon & Accept with expansion & Two 40 kg commodity resupply lots are
manifested before service.\\
Forty spacecraft at 20 kg each within one year & Reject or defer & Aggregate
800 kg demand cannot meet the requested deadline at the required reliability
with the candidate architecture.\\
\bottomrule
\end{tabular}
\end{table}

\section{Discussion}
\label{sec:discussion}

The framework makes the commercial recommendation traceable to launch slots,
deployed assets, inventory flows, and service routes. Its rolling state also
prevents the same capacity from being promised twice.

Several limitations remain. The analytical transfer model does not provide a
flight-ready trajectory. Launch scenarios require calibration and back-testing.
Service duration, failure modes, degradation, ground operations, and recurring
costs require greater fidelity. Exact enumeration is suitable only for small
candidate sets. A full numerical study must compare sequential acceptance
policies, measure the value of launch information, and test sensitivity to
contractual loss boundaries.

Validation should therefore proceed in stages: verify physical feasibility
against trajectory and scheduling tools; back-test uncertainty models on
historical launch calendars; and evaluate acceptance policies using synthetic
and partner-provided contract portfolios.

\section{Conclusions}
\label{sec:conclusion}

This paper presented a rolling-horizon framework for accepting and planning OOS
SLAs. It connects customer contracts to fleet state, inventory, prior
commitments, uncertain launch capacity, infrastructure deployment, and vehicle
routes. Multi-objective architecture search supplies operating concepts, while
manifest selection and scenario evaluation determine whether contractual loss
boundaries can be met. Reliability emerges from the model and includes launch
disruption. Retaining the feasible archive makes counterfactual recommendations
operationally defensible.

Future work will calibrate the joint launch-access model, replace screening
transfers with validated trajectory estimates, scale the manifest problem, and
evaluate the economic value of integrated rolling SLA acceptance.

\section*{Data and code availability}
The present implementation is a research demonstrator. The data and code
availability statement will be completed for submission.

\section*{Declaration of competing interest}
The authors declare no known competing financial interests or personal
relationships that could have appeared to influence this work.

\section*{Acknowledgements}
Acknowledgements will be added for submission.

\begin{thebibliography}{00}
\bibitem[Edelbaum(1961)]{edelbaum1961}
Edelbaum, T.N., 1961. Propulsion requirements for controllable satellites.
ARS Journal 31(8), 1079--1089.

\bibitem[Tricoire(2012)]{tricoire2012}
Tricoire, F., 2012. Multi-directional local search.
Computers \& Operations Research 39(12), 3089--3101.
\end{thebibliography}

\end{document}
